\documentclass[10pt,a4paper]{amsart}
\usepackage{amsmath,amssymb,amsthm,graphicx}
\usepackage{verbatim}
\newenvironment{reference}{\comment}{\endcomment}
%\newenvironment{reference}{}{}
\newenvironment{slogan}{\comment}{\endcomment}
\newenvironment{history}{\comment}{\endcomment}

% For commutative diagrams we use Xy-pic
\usepackage[all]{xy}

% We use 2cell for 2-commutative diagrams.
\xyoption{2cell}
% Theorem environments.
%
\theoremstyle{plain}
\newtheorem{theorem}[subsection]{Theorem}
\newtheorem{proposition}[subsection]{Proposition}
\newtheorem{lemma}[subsection]{Lemma}

\theoremstyle{definition}
\newtheorem{definition}[subsection]{Definition}
\newtheorem{example}[subsection]{Example}
\newtheorem{exercise}[subsection]{Exercise}
\newtheorem{situation}[subsection]{Situation}

\theoremstyle{remark}
\newtheorem{remark}[subsection]{Remark}
\newtheorem{remarks}[subsection]{Remarks}

\numberwithin{equation}{subsection}

% Macros
%
\def\lim{\mathop{\mathrm{lim}}\nolimits}
\def\colim{\mathop{\mathrm{colim}}\nolimits}
\def\Spec{\mathop{\mathrm{Spec}}}
\def\Hom{\mathop{\mathrm{Hom}}\nolimits}
\def\Ext{\mathop{\mathrm{Ext}}\nolimits}
\def\SheafHom{\mathop{\mathcal{H}\!\mathit{om}}\nolimits}
\def\SheafExt{\mathop{\mathcal{E}\!\mathit{xt}}\nolimits}
\def\Sch{\mathit{Sch}}
\def\Mor{\mathop{\mathrm{Mor}}\nolimits}
\def\Ob{\mathop{\mathrm{Ob}}\nolimits}
\def\Sh{\mathop{\mathit{Sh}}\nolimits}
\def\NL{\mathop{N\!L}\nolimits}
\def\CH{\mathop{\mathrm{CH}}\nolimits}
\def\proetale{{pro\text{-}\acute{e}tale}}
\def\etale{{\acute{e}tale}}
\def\QCoh{\mathit{QCoh}}
\def\Ker{\mathop{\mathrm{Ker}}}
\def\Im{\mathop{\mathrm{Im}}}
\def\Coker{\mathop{\mathrm{Coker}}}
\def\Coim{\mathop{\mathrm{Coim}}}

% Boxtimes
%
\DeclareMathSymbol{\boxtimes}{\mathbin}{AMSa}{"02}

%
% Macros for moduli stacks/spaces
%
\def\QCohstack{\mathcal{QC}\!\mathit{oh}}
\def\Cohstack{\mathcal{C}\!\mathit{oh}}
\def\Spacesstack{\mathcal{S}\!\mathit{paces}}
\def\Quotfunctor{\mathrm{Quot}}
\def\Hilbfunctor{\mathrm{Hilb}}
\def\Curvesstack{\mathcal{C}\!\mathit{urves}}
\def\Polarizedstack{\mathcal{P}\!\mathit{olarized}}
\def\Complexesstack{\mathcal{C}\!\mathit{omplexes}}
% \Pic is the operator that assigns to X its picard group, usage \Pic(X)
% \Picardstack_{X/B} denotes the Picard stack of X over B
% \Picardfunctor_{X/B} denotes the Picard functor of X over B
\def\Pic{\mathop{\mathrm{Pic}}\nolimits}
\def\Picardstack{\mathcal{P}\!\mathit{ic}}
\def\Picardfunctor{\mathrm{Pic}}
\def\Deformationcategory{\mathcal{D}\!\mathit{ef}}
\usepackage{cleveref}
\crefname{originaltheorem}{original theorem}{original theorems}
\Crefname{originaltheorem}{Original theorem}{Original theorems}
\makeatletter
\newlabel{theorem-wedderburn}{{0747}{source}{}{sourcepointer-1}{}}
\newlabel{theorem-wedderburn@cref}{{[originaltheorem][0][]0747}{[1][0][]source}{}{}{}}
\makeatother
\crefname{originallemma}{original lemma}{original lemmas}
\Crefname{originallemma}{Original lemma}{Original lemmas}
\makeatletter
\newlabel{lemma-matrix-algebras}{{074D}{source}{}{sourcepointer-2}{}}
\newlabel{lemma-matrix-algebras@cref}{{[originallemma][0][]074D}{[1][0][]source}{}{}{}}
\makeatother
\title{Modules over a finite simple algebra}
\author{The Stacks Project Authors}
\date{stacks-project:brauer.tex:local-snapshot-20261012; GNU FDL 1.2 or later}
\begin{document}
\maketitle
\noindent\textit{Editorial note.} All delivered mathematical text is the author's own, copied from the pinned Stacks Project chapter file sites-cohomology.tex; no statement or proof step was rewritten, reordered or removed.\par\medskip
\noindent\textit{Editorial note.} The proof is the author's complete argument for all six assertions; each is attributed to the named structure theorem, the named equivalence or an explicit description of the module, and no step is left to the reader.\par\medskip
\noindent\textit{Editorial note.} Every labelled reference occurring inside the counted spans is carried in assembly.external\_references as a non-load-bearing navigation pointer (same-chapter ones at their definition site, cross-chapter ones at the invocation inside the counted span), so no load-bearing mathematics is replaced by a stub. No bibliography citation occurs inside the counted spans.\par\medskip
\noindent\textit{Editorial note.} The same-name check for this card scanned every staged Stacks chapter of this repository -- all sources/stacks-*/pinned-original-source.tex files, the card's own chapter excluded -- and the labels found are recorded in dedup.same\_name\_cross\_chapter on the card.\par\medskip
\noindent\textit{Editorial note.} The preamble is the chapter's own retained minimal preamble (sources/stacks-sites-cohomology/retained-preamble.tex), a verbatim excerpt of the native preamble: the Xy-pic setup and the native theorem declarations with the mathematical macros.\par\medskip
\noindent\textit{Editorial note.} The shared gate prints 'occupancy: published-overlap: brauer.tex 483-508' for this candidate while still returning RESULT=OK. That is the known cross-coordinate false overlap of LANE\_BRIEF 10.5: the public item 062 carries upstream RAW line numbers (483-508) while this card carries pinned coordinates, and for brauer the pinned file has an inline preamble of 221 lines, so the raw interval corresponds to pinned lines about 700-725. The true public proposition there is theorem-skolem-noether (Stacks tag 074Q), which this lane excluded by both the gate and the tag intersection; the tag of this card is not the public tag, so this card is not a duplicate.\par\medskip
\noindent\textit{Allowed standard prerequisites (editorial).}\begin{itemize}
\item Wedderburn's theorem. Same chapter, Theorem theorem-wedderburn. Supplies the description of a finite simple algebra as a matrix algebra over a finite skew field extension of the ground field, which reduces everything to modules over a skew field.
\item Equivalence of module categories for a matrix algebra. Same chapter, Lemma lemma-matrix-algebras. Supplies the equivalence between modules over the matrix algebra and modules over the skew field, under which the statements about simple modules, direct sums and dimensions are transported.
\end{itemize}\medskip
\ifdefined\hypertarget\hypertarget{sourcepointer-1}{}\fi
\noindent\textit{Source reference.} Original theorem 0747: sources/stacks-brauer/pinned-original-source.tex, source lines 358--358. Navigation pointer invoked at proof line 519; the counted span contains the pointer only and the cited result is not reproduced inside it, so the pointer is registered as a non-load-bearing navigation stub in the established form for Stacks material. No bibliography citation is involved.\par\medskip
\ifdefined\hypertarget\hypertarget{sourcepointer-2}{}\fi
\noindent\textit{Source reference.} Original lemma 074D: sources/stacks-brauer/pinned-original-source.tex, source lines 477--477. Navigation pointer invoked at proof line 522; the counted span contains the pointer only and the cited result is not reproduced inside it, so the pointer is registered as a non-load-bearing navigation stub in the established form for Stacks material. No bibliography citation is involved.\par\medskip
\begin{lemma}
\label{lemma-simple-module-unique}
Let $A$ be a finite simple $k$-algebra.
\begin{enumerate}
\item There exists exactly one simple $A$-module $M$ up to isomorphism.
\item Any finite $A$-module is a direct sum of copies of a simple module.
\item Two finite $A$-modules are isomorphic if and only if they
have the same dimension over $k$.
\item If $A = \text{Mat}(n \times n, K)$ with $K$ a finite skew field
extension of $k$, then $M = K^{\oplus n}$ is a simple $A$-module and
$\text{End}_A(M) = K^{op}$.
\item If $M$ is a simple $A$-module, then $L = \text{End}_A(M)$
is a skew field finite over $k$ acting on the left on $M$, we have
$A = \text{End}_L(M)$, and the centers of $A$ and $L$ agree.
Also $[A : k] [L : k] = \dim_k(M)^2$.
\item For a finite $A$-module $N$ the algebra $B = \text{End}_A(N)$ is a
matrix algebra over the skew field $L$ of (5). Moreover $\text{End}_B(N) = A$.
\end{enumerate}
\end{lemma}

\begin{proof}
By
Theorem \ref{theorem-wedderburn}
we can write $A = \text{Mat}(n \times n, K)$ for some finite skew
field extension $K$ of $k$. By
Lemma \ref{lemma-matrix-algebras}
the category of modules over $A$ is equivalent to the category of
modules over $K$. Thus (1), (2), and (3) hold
because every module over $K$ is free. Part (4) holds
because the equivalence transforms the $K$-module $K$
to $M = K^{\oplus n}$. Using $M = K^{\oplus n}$ in (5)
we see that $L = K^{op}$. The statement about the center of $L = K^{op}$
follows from
Lemma \ref{lemma-matrix-algebras}.
The statement about $\text{End}_L(M)$ follows from the explicit form
of $M$. The formula of dimensions is clear.
Part (6) follows as $N$ is isomorphic to a direct sum of
copies of a simple module.
\end{proof}

\end{document}
